\chapter{最后要看的作业}


\chapter{第一章作业}
\begin{dxtips}
  证明双射需要，证明端点也就是极限点也能映射到两端，f单调且连续  
\end{dxtips}
\begin{homework}[1.1]
建立从集合 $A$ 到集合 $B$ 的双映射：
\begin{enumerate}
  \item $A = \omega, \quad B = \omega + 1$；
  \item $A = \mathbb{N}, \quad B = \{-1,-2,0\} \cup \mathbb{N}$；
  \item $A = \mathbb{N}, \quad B = \mathbb{Z}$；
  \item $A = (0,2), \quad B = [0,2)$；
  \item $A = (0,2), \quad B = (5,10)$；
  \item $A = \mathbb{R}, \quad B = \left(-\tfrac{\pi}{2}, \tfrac{\pi}{2}\right)$；
  \item $A = \mathbb{R}, \quad B = (0,2)$。
\end{enumerate}
\end{homework}

\begin{homework}[1.2]
建立从集合 $A$ 到集合 $B$ 的满映射：
\begin{enumerate}
  \item $A = \{5,8,9,2,6\}, \quad B = \{a,b,c,d\}$；
  \item $A = \{a,b,e\}, \quad B = \{a,b,c\}$；
  \item $A = \mathbb{Z}, \quad B = \{0,1\}$。
\end{enumerate}
\end{homework}

\begin{homework}[1.3]
证明 $|\mathbb{R}| = |(-\tfrac{\pi}{2}, \tfrac{\pi}{2})|$。
\end{homework}

\begin{homework}[1.14]
证明有理数集 $\mathbb{Q}$ 是可数集。
\end{homework}

\begin{homework}[1.15]
证明 $B=\{(a,b): a\in \mathbb{Q}, b\in \mathbb{Q}\}$ 是可数集，其中 $\mathbb{Q}$ 是有理数集，$(a,b)$ 是开区间。
\end{homework}
\begin{dxtips}
    \mathbb{Q}\times\mathbb{Q}=\{\langle a,b\rangle:a\in\mathbb{Q},\ b\in\mathbb{Q}\}. \quad (a,b)=\{x\in\mathbb{R}:a<x<b\}.
\end{dxtips}
\begin{homework}[1.16]
证明 $B=\{(a,b)\times (c,d): a<b<c<d,\, a\in \mathbb{Q}, b\in \mathbb{Q}, c\in \mathbb{Q}, d\in \mathbb{Q}\}$ 是可数集，其中 $\mathbb{Q}$ 是有理数集，$(a,b)$ 是开区间。
\end{homework}

\begin{homework}[1.17]
写出下列集合 $X$ 的幂集 $P(X)$，并写出 $|X|$ 与 $|P(X)|$：
\begin{enumerate}[(1)]
  \item $X = \varnothing$；
  \item $X = \{\varnothing, \{\varnothing\}\}$；
  \item $X = \{a,b,c\}$。
\end{enumerate}
\end{homework}
\begin{dxtips}
    是的，$X=\prod_{i\in\mathbb{N}}X_i$ 且 $X_i=\{0,1\}$ 时，有 $X=\{0,1\}^{\mathbb{N}}$，它由所有 $0$--$1$ 的无限序列组成，与 $\mathcal{P}(\mathbb{N})$ 等势，因此基数为 $2^{\aleph_0}$，是不可数集（亦可由 Cantor 对角线法证明）。
\end{dxtips}
\begin{homework}[1.18]
下列集合是可数集还是不可数集？若是可数集，说明是否有限集。
\begin{enumerate}[(1)]
  \item $X = \mathbb{Z}$；
  \item $X = \mathbb{R}$；
  \item $X = \mathbb{R}\cap \mathbb{Q}$；
  \item $X = \mathbb{Q}\times \mathbb{Z}$；
  \item $X = \prod_{i\in\mathbb{N}} X_i, \ X_i=\{0,1\},\ i\in\mathbb{N}$；
  \item $X = \prod_{i=1}^{10} X_i, \ X_i=\{0,1\},\ i\leq 10$；
  \item $X = \prod_{i\in\mathbb{N}} X_i, \ X_i=\{0,1\},\ i\leq 10 \text{时},\ X_i=\{0,1\},\ i\geq 11 \text{时},\ X_i=\{1\}$；
  \item $X = \prod_{i\in\mathbb{N}} X_i, \ X_i=\{0\},\ i\in\mathbb{N}$；
  \item $X = \prod_{i=1}^{10} X_i, \ X_i=\mathbb{Q},\ i\leq 10$；
  \item $X = \prod_{i=1}^{\infty} X_i, \ X_i=\mathbb{Q},\ i\leq 10$。
\end{enumerate}
\end{homework}
\begin{dxtips}
    互相包含，建立映射把右边分成两半方便找和有理数对应。
\end{dxtips}
\begin{homework}[1.19]
证明 $|(0,1)| = |(0,1)\cap I|$。
\end{homework}
\chapter{2}

\subsection*{第三周作业开集闭集判断}

\begin{exercise}
设 $\mathbb{R}$ 是实数直线，拓扑是通常拓扑，判断下列集合 $A$ 是否是开集：  
\begin{enumerate}
  \item $A = (1,2);$
  \item $A = (1,2) \cup \{3\};$
  \item $A = \mathbb{Q}, \quad \mathbb{Q} \text{ 为有理数集};$
  \item $A = [3,4);$
  \item $A = \mathbb{R}\setminus\left\{\tfrac{1}{n} : n\in\mathbb{N}\right\};$
  \item $A = \mathbb{R}\setminus\mathbb{Z};$
  \item $A = (1,3]\cup(2,4).$
\end{enumerate}
\end{exercise}

\begin{exercise}
证明：空间 $X$ 是离散空间当且仅当 $X$ 中的每个单点集是开集。
\end{exercise}
\begin{dxtips}
    关键：$\{1/n:n\in\mathbb{N}\}$ 不是开集（通常拓扑下）。因为对任意 $1/n$，不管你取多小的开区间 $(1/n-\varepsilon,\,1/n+\varepsilon)$，里面都会有很多不是 $1/k$ 的实数点，所以这个开区间不可能被包含在 $\{1/n\}$ 这种“离散点集”里，因此 $1/n$ 不是内点，整个集合没有内点，就不是开集。
    并不是不是开集就是闭集，而是开集的补是闭，闭的补是开。
\end{dxtips}
\begin{exercise}
设 $\mathbb{R}$ 是实数直线，拓扑是通常拓扑，判断下列集合 $A$ 是否是闭集：  
\begin{enumerate}
  \item $A = [1,2];$
  \item $A = [1,2]\cup\{3\};$
  \item $A = \mathbb{Q};$
  \item $A = [3,4];$
  \item $A = \mathbb{R}\setminus\left\{\tfrac{1}{n}:n\in\mathbb{N}\right\};$
  \item $A = \mathbb{Z};$
  \item $A = [1,3)\cup(2,4];$
  \item $A = \left\{\tfrac{1}{n}:n\in\mathbb{N}\right\}.$
\end{enumerate}
\end{exercise}

\section*{第四周作业，三种空间}
\begin{dxtips}
    第二可数空间需要有一个可数开邻域基B，任意开集都有Bx属于他也就是比他小。这个可数开邻域基就选（q-1/n，q+1/n），任意x属于U存在ε大于0（x- ε，x+ε）属于U然后让1/n比ε小即可
\end{dxtips}
\begin{dxtips}
    林德洛夫任意开覆盖都要可数子覆盖，离散空间就取单点集做开覆盖，有限个只能覆盖有限个单点集。
\end{dxtips}
\begin{dxtips}
    基要求可以生成任何开集，并起来能组成整个空间，任何开集都可以找到Bx属于B，Bx属于开集
\end{dxtips}
\begin{homework}
[4]
\begin{enumerate}
  \item 证明 $\mathbb{R}$ 在通常拓扑下是第二可数空间；
  \item 证明：如果 $X$ 是不可数的离散拓扑空间，则 $X$ 不是 Lindelöf 空间；
  \item 如果 $X$ 是离散拓扑空间，则 $\mathcal{B} = \{\{x\}: x \in X\}$ 是空间 $X$ 的一个基。
\end{enumerate}
\end{homework}
\begin{dxtips}
    第一可数空间需要每一个x都有可数基，这个基就取（x-1/n，x+1/n）然后还是x属于U，（x-ε，x+ε）属于U然后1/n小于ε即可
\end{dxtips}
\begin{homework}
[5] 
证明：$\mathbb{R}$ 在通常拓扑下是第一可数空间
\end{homework}


\section{第五周作业-闭包内点}
\begin{dxtips}
    闭包是加入极限点，导集是闭包去掉孤立点，内点集是所有开集的并是开集
\end{dxtips}
\begin{homework}[6]
设 $\mathbb{R}$ 是实数直线，拓扑为通常拓扑。对 $A \subset \mathbb{R}$，分别写出其闭包 $\overline{A}$、导集 $A^d$、内点集 $A^\circ$。

\begin{enumerate}[(1)]
  \item $A = (1,2)$；
  \item $A = (1,2)\cup\{3\}$；
  \item $A = \mathbb{Q}$；
  \item $A = [3,4)$；
  \item $A = \left\{\tfrac{1}{n} : n\in\mathbb{N}\right\}$；
  \item $A = \{\pi\}$；
  \item $A = \mathbb{Z}$；
  \item $A = \left\{2 - \tfrac{1}{n} : n\in\mathbb{N}\right\}$；
  \item $A = \mathbb{I}$（无理数集）。
\end{enumerate}
\end{homework}

\begin{homework}[7]
设 $\mathbb{R}^2$ 是平面上的实数直线空间，拓扑为通常拓扑。对 $A \subset \mathbb{R}^2$，分别写出其闭包 $\overline{A}$、导集 $A^d$、内点集 $A^\circ$。

\begin{enumerate}[(1)]
  \item $A = (1,2)\times(3,4)$；
  \item $A = (1,2)\times\{3\}$；
  \item $A = \mathbb{Q}\times\mathbb{Q}$；
  \item $A = \{(x,y): x^{2}+y^{2}<4\}$；
  \item $A = \{(x,y): x+y=4\}$；
  \item $A = \mathbb{R}\times\{0\}$；
  \item $A = \mathbb{R}^{2}\setminus\{(x,y):x^{2}+y^{2}=4\}$。
\end{enumerate}
\end{homework}

\begin{homework}[14]
证明 $\mathbb{R}$ 在通常拓扑下是可分的。
\end{homework}
\section{各种空间}

\begin{homework}[第 7 题]
设 $\mathbb{R}$ 是实数直线，拓扑为通常拓扑。  
$A \subset \mathbb{R}^2$，分别写出 $\overline{A}$ 与 $A^\circ$。

\begin{enumerate}[(1)]
  \item $A = (1,2) \times (3,4);$
  \item $A = (1,2) \times \{3\};$
  \item $A = \mathbb{Q} \times \mathbb{Q};$
  \item $A = \{(x,y): x^2 + y^2 < 4\};$
  \item $A = \{(x,y): x + y = 4\};$
  \item $A = \mathbb{R} \times \{0\};$
  \item $A = \mathbb{R}^2 \setminus \{(x,y): x^2 + y^2 = 4\}.$
\end{enumerate}
\end{homework}

\begin{homework}[第 8 题]
证明：$\mathbb{R}^2$ 是可分空间。
\end{homework}

\begin{homework}[第 9 题]
证明 $\mathbb{R}^2$ 是第一可数空间。
\end{homework}
\begin{dxtips}
    注意题目给你什么如果和开覆盖，开集相关的条件但是给你的手闭子空间，此时就要学会取闭子空间的补集做文章了。这题就巧妙利用了x\Y还是开集所以可以和Dj∩Y=Uj，也就是X中取的开集与Y交得到Uj的开集组成开覆盖。这个写法可以温习一下
\end{dxtips}
\begin{homework}[10]
证明 Lindelöf 空间的闭任意子空间是 Lindelöf 空间。
\end{homework}
\begin{dxtips}
    为什么要闭子空间就是要他的补集是开集，然后先取子空间的覆盖，子空间的覆盖中每一个开集都可以写成Y和X中开集的交也就是VU们，VU们和X/Y组成的开覆盖就可以取有限的子覆盖了就把Vu从无限变成可数个了。他Vu的U也变成可数个了因为是一一对应的。
\end{dxtips}

\begin{proof}
设 $X$ 为 Lindel\"of 空间，$Y\subset X$ 为闭子集，且 $Y$ 带子空间拓扑。
取 $\mathcal U$ 为 $Y$ 的任意一族开集，使其覆盖 $Y$。

由子空间拓扑的定义，对每个 $U\in\mathcal U$，存在 $X$ 中的开集 $V_U$，
使得
\[
U = Y \cap V_U.
\]
由于 $Y$ 是 $X$ 的闭子集，故 $X\setminus Y$ 是 $X$ 中的开集。
于是
\[
\{\,V_U : U\in\mathcal U\,\}\cup\{X\setminus Y\}
\]
构成 $X$ 的一个开覆盖。

因为 $X$ 是 Lindel\"of 空间，存在可数子族
$\{V_{U_n}\}_{n\in\mathbb N}$，使得
\[
X \subset (X\setminus Y)\ \cup\ \bigcup_{n=1}^{\infty} V_{U_n}.
\]
两边同时与 $Y$ 取交，并注意到 $Y\cap(X\setminus Y)=\varnothing$，得到
\[
Y \subset \bigcup_{n=1}^{\infty} (Y\cap V_{U_n})
= \bigcup_{n=1}^{\infty} U_n.
\]
因此 $\{U_n\}_{n\in\mathbb N}$ 是 $\mathcal U$ 的一个可数子覆盖。

由 $\mathcal U$ 的任意性可知，$Y$ 的任意开覆盖都有可数子覆盖，
故 $Y$ 是 Lindel\"of 空间。
\end{proof}
\begin{dxtips}
    开子空间就是为了和稠密的Q相交不是空集。
\end{dxtips}
\begin{homework}[11]
设 $\mathbb{R}$ 是实数直线，拓扑是通常拓扑。证明 $\mathbb{R}$ 的每个开子空间都是可分空间。
\end{homework}

\begin{homework}[12]
证明积空间 $\mathbb{R}^2$ 是第二可数空间。
\end{homework}
\chapter{3}
\section{第七周作业,连续映射}
\begin{dxtips}
    连续映射保证了开集的原像还是开集，但是保证的是y不是Y的子空间所以子空间的开集要写成子空间交Y的一个开集的形式，显然子空间这块就有可能破坏性质，所以我们一般取的这个子空间都是完整的像集这样保证逆映射回去就是X肯定是开集，然后再交集还是开集就找到我们想要的子空间中开集的原像开集了
\end{dxtips}
\begin{example}
    是不是比如说R连续映射到了[0,1),在[0,1）这个子空间中[0,1)就是开集所以不能直接用子空间中的开集
\end{example}
    

\begin{homework}[2.1]
可分空间的每个开子空间也都是可分空间。
\end{homework}

\begin{homework}[13]
设 $\mathbb{R}$ 是实数直线，拓扑是通常拓扑，$A\subset \mathbb{R}^2$。  
判断下列集合是否是积空间 $\mathbb{R}^2$ 中的开集？是否是闭集？
\begin{enumerate}[(1)]
  \item $A = (1,2)\times(3,4)$；
  \item $A = (1,2)\times\{3\}$；
  \item $A = \mathbb{Q}\times\mathbb{Q}$；
  \item $A = \{(x,y): x^2+y^2<4\}$；
  \item $A = \{(x,y): x+y=4\}$；
  \item $A = \mathbb{R}\times\{0\}$；
  \item $\mathbb{R}^2\setminus\{(x,y):x^2+y^2=4\}$。
\end{enumerate}
\end{homework}
\begin{dxtips}
    第一可数需要子空间Y中每一个点y都有可数邻域基，那y现在X中找可数邻域基再和Y做交集找到Y中的可数邻域基。第二可数同理整个X有可数开邻域基，任意子空
\end{dxtips}
\begin{homework}[2.2]
每个第一可数（第二可数）空间的子空间也是第一可数（第二可数）空间。
\end{homework}
\begin{proof}
\textbf{（第一可数）}
设 $(X,\mathcal T)$ 是第一可数空间，$Y\subset X$，$y\in Y$。
由 $X$ 的第一可数性，存在 $y$ 在 $X$ 中的一个可数邻域基
\[
\{U_n:n\in\mathbb N\},
\]
其中每个 $U_n$ 是 $X$ 中的开集，且对任意 $y$ 的开邻域 $U\subset X$，
存在 $n$ 使得 $U_n\subset U$。

由于 $Y$ 取子空间拓扑，$Y$ 中含 $y$ 的任一开邻域 $V$
都可表示为
\[
V=U\cap Y,
\]
其中 $U$ 是 $X$ 中含 $y$ 的开集。
由上述性质，存在 $n$ 使得 $U_n\subset U$，于是
\[
U_n\cap Y\subset U\cap Y=V.
\]
并且 $U_n\cap Y$ 是 $Y$ 中的开集。
因此
\[
\{U_n\cap Y:n\in\mathbb N\}
\]
构成 $y$ 在子空间 $Y$ 中的一个可数邻域基。
由 $y$ 的任意性，$Y$ 是第一可数空间。

\medskip
\textbf{（第二可数）}
若 $(X,\mathcal T)$ 是第二可数空间，设 $\mathcal B$ 是 $X$ 的一个可数拓扑基。
定义
\[
\mathcal B_Y=\{B\cap Y: B\in\mathcal B\}.
\]
显然 $\mathcal B_Y$ 仍是可数集。

下面验证 $\mathcal B_Y$ 是 $Y$ 的拓扑基。
设 $y\in Y$，$V$ 是 $Y$ 中含 $y$ 的开集，
则存在 $X$ 中的开集 $U$ 使
\[
V=U\cap Y.
\]
由于 $\mathcal B$ 是 $X$ 的基，存在 $B\in\mathcal B$ 使得
\[
y\in B\subset U.
\]
于是
\[
y\in B\cap Y\subset U\cap Y=V,
\]
且 $B\cap Y\in\mathcal B_Y$。
因此 $\mathcal B_Y$ 是 $Y$ 的拓扑基。

故 $Y$ 也是第二可数空间。
\end{proof}
\begin{dxtips}
离散拓扑里面所有集合都是开集自然无论逆映射回来是什么都是开集，然后平凡拓扑只有空集和自己本身所以，空集逆映射还是空集是开集，Y逆映射就是X
\end{dxtips}

\begin{homework}[1]
证明：\(X\) 与 \(Y\) 是拓扑空间且 \(X\) 是离散拓扑空间。  
则任一映射 \(f:X\to Y\) 都是连续映射。
\end{homework}

\begin{homework}[2]
证明：\(X\) 与 \(Y\) 是拓扑空间且 \(Y\) 是平凡拓扑空间。  
则任一映射 \(f:X\to Y\) 都是连续映射。
\end{homework}
\begin{dxtips}
    也就说只用找f（x）映射到（1，3）的原像就行，没说原像映射过去需要充满（1，3）
\end{dxtips}
\begin{dxtips}
    是的，$(0,+\infty)$ 在 $\mathbb{R}$ 的通常拓扑下是开集。
    所以得是[0，正无穷]
\end{dxtips}
\begin{homework}[5]
设 \(\mathbb{R}\) 是实数直线，拓扑为通常拓扑。  
定义映射 \(f:\mathbb{R}\to\mathbb{R}\) 为
\[
f(x)=
\begin{cases}
2, & x\ge0,\\
5, & x<0.
\end{cases}
\]
证明：\(f\) 不是连续映射。
\end{homework}

\begin{homework}[6]
设 \(\mathbb{R}\) 是实数直线，拓扑为离散拓扑。  
定义映射 \(f:\mathbb{R}\to\mathbb{R}\) 为
\[
f(x)=
\begin{cases}
2, & x\ge0,\\
5, & x<0.
\end{cases}
\]
证明：\(f\) 是连续映射。
\end{homework}

\section{第八周作业，连续映射}
\begin{dxtips}
    映射过去等于自身，逆映射也等于自身。拓扑里面的都是开集所以可以证明
\end{dxtips}
\begin{dxtips}
    别忘了X本身
\end{dxtips}
\begin{homework}[7]
设 $X=\{a,b,c\}$，拓扑为 
\[
\mathcal{T}_X=\{X,\varnothing,\{a\},\{a,b\},\{a,c\}\},
\]
又设
\[
Y=\{a,b,c\},\quad 
\mathcal{T}_Y=\{Y,\varnothing,\{a\},\{a,b\},\{a,c\}\}.
\]
若映射 $f:X\to Y$ 满足 $f(x)=x$，问：$f$ 是否连续？
\end{homework}

\begin{homework}[8]
设 $X=\{a,b,c\}$，拓扑为 
\[
\mathcal{T}_X=\{X,\varnothing,\{a\},\{a,b\},\{a,c\}\},
\]
又设
\[
Y=\{a,b\},\quad 
\mathcal{T}_Y=\{Y,\varnothing,\{b\}\}.
\]
若映射 $f:X\to Y$ 满足 
\[
f(a)=b,\quad f(b)=b,\quad f(c)=a,
\]
问：$f$ 是否连续？
\end{homework}

\begin{homework}[9]
设 $X=\{a,b,c\}$，拓扑为 
\[
\mathcal{T}_X=\{X,\varnothing,\{a\},\{a,b\},\{a,c\}\},
\]
又设
\[
Y=\{a,b\},\quad 
\mathcal{T}_Y=\{Y,\varnothing,\{b\}\}.
\]
若映射 $f:X\to Y$ 满足 
\[
f(a)=a,\quad f(b)=b,\quad f(c)=a,
\]
问：$f$ 是否连续？
\end{homework}
\begin{dxtips}
    X是整个像集，只要逆映射就能逆回整个集合所以不管对他交逆映射回去只会剩下他的另一部分
\end{dxtips}
\begin{dxtips}
    如果 $X$ 到 $Y$ 是连续映射也就说证明充分性，
$Y$ 中的开集都逆映射是开集了。
$f(X)$ 只是 $Y$ 中的子集，
\textcolor{red}{但 $f(X)$ 中的开集不一定是 $Y$ 中的开集}，
而是可以写成 $Y$ 中某个开集 $V$ 与 $f(X)$ 的交，
即 $U=V\cap f(X)$，
因此自然逆映射过去仍然是开集。

必要性方向中，因为 $X$ 到 $f(X)$ 是连续映射，
$f(X)$ 中所有的开集 $U$ 都能写成
$Y$ 中开集 $V$ 和 $f(X)$ 的交的形式，
这也是子集子空间中常用的表示方式。
由于 $U$ 的逆映射是开集，
而 $U$ 的逆映射等于 $V$ 的逆映射与整个 $X$ 的交，
仍然是 $V$ 的逆映射，
所以所有 $Y$ 中开集 $V$ 的逆映射都是开集。
\end{dxtips}
\begin{homework}[12]
证明：映射 $f:X\to Y$ 连续，当且仅当 $f:X\to f(X)$ 连续。
\end{homework}
\begin{dxtips}
    连续的像就是f（x）,没有Y剩下不是像的部分
\end{dxtips}
\begin{homework}[3]
证明：Lindelöf 拓扑空间的连续像是 Lindelöf 空间。
\end{homework}
\begin{dxtips}
    这章重点就是他的像集逆回去会消失，整个映射过去能充满。
\end{dxtips}
\begin{homework}[4]
证明：可分拓扑空间的连续像是可分空间。
\end{homework}
\begin{dxtips}
    f｜A和f不一样他的逆的像本来就要和A做交集。
\end{dxtips}
\begin{homework}[14]
如果 $f:X\to Y$ 是连续映射，$A\subset X$，证明：$f|_A:A\to Y$ 也是连续映射。
\end{homework}
\begin{dxtips}
 现在吧f1，f2，用F投影到f1（x），f2（x）再用把F投影到f1和F投影到f2，因为投影本身就是连续的加上f1，f2都是连续的所以F连续，然后弄出a进行我们要的运算，a还是连续的a和F复合出来就是f1+f2连续。F连续哪块也能直接写因为分量都连续所以连续
\end{dxtips}
\begin{homework}[3.2]
若 $f_i:X\to\mathbb{R}$ 是连续映射，$i\in\{1,2\}$，证明：$f_1+f_2$、$f_1-f_2$、$f_1f_2$ 及 $f_1,f_2$ 都是连续映射。
\end{homework}
\begin{proof}
设 $f_1,f_2:X\to\mathbb{R}$ 连续，定义
\[
F:X\to\mathbb{R}^2,\qquad F(x)=(f_1(x),f_2(x)).
\]
记 $\pi_1,\pi_2:\mathbb{R}^2\to\mathbb{R}$ 为两个坐标投影，则 $\pi_1,\pi_2$ 连续且
\[
\pi_1\circ F=f_1,\qquad \pi_2\circ F=f_2.
\]
由分量连续可知 $F$ 连续。再定义加法映射
\[
a:\mathbb{R}^2\to\mathbb{R},\qquad a(u,v)=u+v,
\]
其为连续映射（例如由欧氏空间中多项式函数连续可知）。于是
\[
f_1+f_2=a\circ F,
\]
为连续映射（连续映射的复合仍连续）。因此 $f_1+f_2$ 连续。
\end{proof}
\begin{homework}[2.1]
证明：可分空间的每个开子空间也是可分空间。
\end{homework}

\section{第九周作业，闭映射}

\begin{homework}[13]
证明函数 $y = f(x) = \sin x \ (x \in \mathbb{R})$ 不是从 $\mathbb{R}$ 到 $[-1,1]$ 的闭映射。
\end{homework}
\begin{dxtips}
    这里+2nπ就是为了保证函数值可控的同时不收敛，就不用加入极限点。n选大于1也是不让sin0=0或者sinπ=0出现。是闭集的原因就是自己等于自己的闭包，压根没有极限点了。
\end{dxtips}
\begin{proof}
令
\[
A=\left\{\frac{\pi}{n}+2n\pi:\ n\in\mathbb{N},\ n>1\right\},
\]
则 $A$ 是 $\mathbb{R}$ 中的闭集。而
\[
f(A)=\left\{\sin\frac{\pi}{n}:\ n\in\mathbb{N},\ n>1\right\}.
\]
由于
\[
\lim_{n\to\infty}\sin\frac{\pi}{n}=0,
\]
因此
\[
\overline{f(A)}=f(A)\cup\{0\},
\qquad
0\notin f(A).
\]
从而 $f(A)$ 不是 $[-1,1]$ 中的闭集。故
\[
y=\sin x\ (x\in\mathbb{R})
\]
不是从 $\mathbb{R}$ 到 $[-1,1]$ 的闭映射。
\end{proof}
\begin{dxtips}
    又是两个包含导致相等。
\end{dxtips}
\begin{homework}[3.6]
设 $f : X \to Y$ 是连续映射，证明：$f$ 是闭映射当且仅当对每个 $A \subset X$ 都有
\[
f(\overline{A}) = \overline{f(A)}.
\]
\end{homework}
\begin{dxtips}
    本身就是定理
    对任意 $\varepsilon>0$，存在点 $x_0$ 的某个邻域 $O(x_0)$，使得只要
$x\in O(x_0)$，就有
\[
|f(x)-f(x_0)|<\varepsilon.
\]
在度量空间（如 $\mathbb{R}$）中，可以取
\[
O(x_0)=\{x:\ |x-x_0|<\delta\},
\]
其中 $\delta>0$，这正是通常的 $\varepsilon$--$\delta$ 定义中的条件 $|x-x_0|<\delta$。
这里还是用的一手拆分。后半部分用分布收敛到0，前半部分用连续的性质加上对x对x0距离也就说δ做限制然后取最小的δ收敛到ε。用集合的方式写就是交集然后所有的点都在附近也就交集不是空
\end{dxtips}
\begin{dxtips}
    还是要证明f（x）-f（x0）小于ε在x和x0一定条件下成立，先把这个拆成两个部分一部分因为本身分母很大加和趋紧于0，另一部分取每一个fn保证x-x0
小于δ f（x0）-f（x0小于ε范围点交集。然后两个相加小于ε就可以了
\end{dxtips}
\begin{homework}[3.3]
若 $f_n : X \to [a,b]$ 是连续映射，$n \in \mathbb{N}$，证明若
\[
f(x) = \sum_{n=1}^{\infty} \frac{f_n(x)}{2^n},
\]
则 $f : X \to [a,b]$ 是连续映射。
\end{homework}
\begin{proof}
对任意 $n\in\mathbb{N}$，由于 $0\le f_n(x)\le 1$，因此
\[
\left|\frac{f_n(x)}{2^n}\right|\le \frac{1}{2^n}.
\]
这样由引理 3.9 可知
\[
\sum_{n=1}^{\infty}\frac{f_n(x)}{2^n}
\]
在 $X$ 上一致收敛到 $f(x)$。

对任意 $x_0\in X$，有
\[
f(x)-f(x_0)=\sum_{n=1}^{\infty}\frac{f_n(x)-f_n(x_0)}{2^n}.
\]
于是对任意 $\varepsilon>0$，存在 $m\in\mathbb{N}$，当 $n\ge m+1$ 时有
\[
\left|\sum_{n=m+1}^{\infty}\frac{f_n(x)-f_n(x_0)}{2^n}\right|
\le \sum_{n=m+1}^{\infty}\frac{1}{2^{n-1}}
= \frac{1}{2^{m-1}}<\frac{\varepsilon}{2}.
\]

对 $n\le m$，由于 $f_n(x)$ 在点 $x_0$ 连续，因此存在 $X$ 中含点 $x_0$ 的开集
$V_n(x_0)$，使得当 $x\in V_n(x_0)$ 时，有
\[
|f_n(x)-f_n(x_0)|<\frac{\varepsilon}{2},\qquad n\le m.
\]

于是令
\[
O(x_0)=\bigcap_{n=1}^{m}V_n(x_0),
\]
则 $O(x_0)$ 是点 $x_0$ 的开邻域，且当 $x\in O(x_0)$ 时，
\[
\left|\sum_{n=1}^{m}\frac{f_n(x)-f_n(x_0)}{2^n}\right|
\le \frac12\left(1-\frac{1}{2^m}\right)\frac{\varepsilon}{1-\frac12}
<\frac{\varepsilon}{2}.
\]

因此当 $x\in O(x_0)$ 时，
\[
|f(x)-f(x_0)|
\le \frac{\varepsilon}{2}+\frac{\varepsilon}{2}
=\varepsilon.
\]
从而 $f:X\to[0,1]$ 连续。
\end{proof}
\begin{homework}[10] 
设 $\mathbb{R}$ 是实数直线，拓扑是通常拓扑，$(-\frac{\pi}{2}, \frac{\pi}{2})$ 是 $\mathbb{R}$ 的子空间。证明 $\mathbb{R}$ 与 $ \left( -\frac{\pi}{2}, \frac{\pi}{2} \right)$ 同胚。
\end{homework}

\begin{homework}[11]
设 $\mathbb{R}$ 是实数直线，拓扑是通常拓扑，$(0,1)$ 与 $(0,2)$ 是 $\mathbb{R}$ 的子空间。证明 $(0,1)$ 与 $(0,2)$ 同胚。
\end{homework}
\chapter{4}
\section{第十周作业，连通空间}

\begin{homework}[1]
\begin{enumerate}
  \item 设 $X=\{a,b,c\}$，$\mathcal{T}=\{X,\varnothing,\{a\},\{a,b\},\{a,c\}\}$。
  \item 设 $Y=\{a,b,c\}$，$\mathcal{T}=\{Y,\varnothing,\{a\},\{a,c\}\}$。
  \item 设 $X=\{a,b,c\}$，$\mathcal{T}=\{X,\varnothing,\{a\},\{b\},\{a,b\},\{b,c\}\}$。
  \item 设 $X=\{a,b,c\}$，$\mathcal{T}=\{X,\varnothing,\{b\},\{a,c\}\}$。
  \item 设 $X=\{a,b,c,d\}$，$\mathcal{T}=\{X,\varnothing,\{b\},\{a,c\},\{d\},\{b,d\},\{a,b,c\}\}$。
\end{enumerate}
\end{homework}

\begin{homework}[2]
证明连通空间的连续像是连通空间。
\end{homework}
\begin{dxtips}
    如果不是区间，（a，b）中有点c不在E上那E做和（-∞，c）和c（正无穷）的交就得到了两个两个相交起来是空集但是并起来等于E的集合E就不连通了了。区间就是区间可以开可以闭可以一边开一边闭。
\end{dxtips}
\begin{homework}[3]
证明 $\mathbb{R}$ 上的连通集都是单点集或是区间。
\end{homework}
\begin{dxtips}
   一族联通集合如果有公共点也就是交集不是空集他们的并就是联通的 
\end{dxtips}
\begin{dxtips}
    包含x的所有联通集的并是，x联通集的最大并，只要包含x的联通集都是他的子集。要证明C的闭包里面的点都是C里面的点，先取y点，C并{y}形成新的集合D，然后把D划分成包含y点开集合u和不包含y点开集合v，这个不包含y的集合肯定是C的一个子集。
\end{dxtips}
\begin{homework}[4.1]
证明空间 $X$ 中包含给定点 $x$ 的所有连通集的并是闭的连通集。
\end{homework}
\begin{dxtips}
    如果B不属于A
    把B分成 B∩A和B交{X/A}，因为A是X中开集所以B交A是开集，因为A还可以是闭集导致X/A还是开集，B交完了还是开集。所以B分成了两个互不相交的开集所以B是联通矛盾所以B是A的子集让B交{X/A}变成空集
\end{dxtips}
\begin{dxtips}
    要给所有连通集的并也就A要找到两个互不相交的开集
    。x属于A的闭包，然后D=x并A把D如果不连通把D分成一个包含x的开集一个不包含X的开集，不包含x的开集显然是A的子集，用V交A=V。包含X的U也要和A交集，怎么证明不是空的？U可以写成D∩G G为X中开集，G是包含x的开集所有根据闭包性质闭包中的点的开集都和原集有交集，闭包操作时不会加入孤立的点。所以A交u就变成A交（D交G）分开交就不是空集了。然后D就是连通的D就属于Ax也就是属于A
\end{dxtips}
\begin{homework}[4.4]
设 $A$ 是 $X$ 中既开又闭的集合，$B$ 是连通集。若 $A\cap B\neq\varnothing$，证明 $B\subset A$。
\end{homework}

\section{第十一周作业，紧集}

\begin{homework}[4.5]
$\,\mathbb R$ 是实数直线，拓扑为通常拓扑。判断下列集合是否连通：
\begin{enumerate}
    \item $Y=(1,2)\cup(2,5)$;
    \item $Y=(1,3)\cup(3,5)$;
    \item $Y=[1,3)\cup(3,5]$;
    \item $Y=[1,3)\cup[3,5]$;
    \item $Y=[1,2]\cup(3,5]$;
    \item $I_Y=\mathbb Q$;
    \item $Y=\{\pi,\pi+1\}$;
    \item $Y=\mathbb Q\cup I$;
    \item $Y=\overline{\mathbb Q}$.
\end{enumerate}
\end{homework}
\begin{dxtips}
    连通+连续项也是连通的，R里面连通的只有单点集和区间，显然像就是区间。0总是f（x）到f（y）里面的所以0在区间内，
\end{dxtips}
\begin{homework}[4.6]
设 $X$ 是连通空间，$f:X\to\mathbb R$ 连续。若 $x,y\in X$ 且 $f(x)f(y)<0$，证明存在 $z\in X$ 使得 $f(z)=0$。
\end{homework}
\begin{dxtips}
    R中去掉0点不连通了，但是R²去掉一个点还是道路连通的，同理R去掉一个点不连通了[1,2)]去掉1还是连通的。
\end{dxtips}
\begin{homework}[4.7]
证明 $\mathbb{R}^2$ 与 $\mathbb{R}$ 不同胚。
\end{homework}

\begin{homework}[4.8]
证明 $\mathbb{R}$ 与区间 $[1,2)$ 不同胚。
\end{homework}

\begin{homework}[4.9]
设 $\mathbb{R}$ 为实数直线，拓扑为通常拓扑，$A\subset\mathbb{R}^2$。判断下列集合是否连通：
\begin{enumerate}
  \item $A=(1,2)\times(3,4)$；
  \item $A=(1,2)\times\{3\}$；
  \item $A=\mathbb{Q}\times\mathbb{Q}$；
  \item $A=\{(x,y):x^{2}+y^{2}<4\}$；
  \item $A=\{(x,y):x+y=4\}$；
  \item $A=\{(x,y):x^{2}+y^{2}<4\}\cup\{(x,y):x^{2}+y^{2}>4\}$。
\end{enumerate}
\end{homework}

\section{第十二周作业}
\begin{definition}[序列的拓扑收敛]
设 $(X,\mathcal T)$ 为拓扑空间，$(x_n)$ 为 $X$ 中的一个序列，$x\in X$。
称
\[
x_n \to x,
\]
当且仅当：对任意包含 $x$ 的开集 $U\in\mathcal T$，存在 $N\in\mathbb N$，使得对所有 $n\ge N$ 都有
\[
x_n\in U.
\]
\end{definition}

\noindent
\textbf{直观表述：} 从某一项开始，序列始终落在任意给定的“$x$ 的开邻域”中。

\medskip
\noindent
\textbf{等价的邻域表述：}
\[
x_n\to x
\quad\Longleftrightarrow\quad
\forall\,\text{$x$ 的邻域 }N_x,\ \exists\,N\in\mathbb N,\ \forall\,n\ge N,\ x_n\in N_x.
\]
\begin{homework}[5.1]
证明在 $T_2$ 拓扑空间中，每个收敛序列有唯一的极限点。
\end{homework}

\begin{homework}[5.2]
在空间 $X$ 中，若序列 $\{x_n\}$ 收敛于点 $x_0$，证明集合
\[
\{x_0\}\cup\{x_n:n\in\mathbb{N}\}
\]
是紧集。
\end{homework}

\begin{homework}[2]
证明紧子空间的连续像都是紧子空间。
\end{homework}
\begin{dxtips}
   直接证明闭集不好整就整明补集是开集。这题就要通过 $T_2$ 空间性质找补集，然后再用有限覆盖找到有限交。$x$ 得在 $X\setminus A$ 中取保证每一个 $x$ 都可以属于一个开集，每一个 $x$ 都有开邻域整个 $X\setminus A$。构造开集的时候可以都用Uy，Vy方便说明对应性。证明U与A没有交集只需要先证明U和Vyi的交集都是空集，也就是U交Vyi属于Uyi交Vyi
\end{dxtips}
\begin{homework}[3]
证明：$T_2$ 空间中的紧集是闭集。
\end{homework}
\begin{dxtips}
    存在很好证明的
\end{dxtips}
\begin{homework}[6]
证明：$T_2$ 空间的子空间是 $T_2$ 空间。
\end{homework}
\begin{dxtips}
    
    
   若存在 $Z\in U\cap V$，则
\[
d(P,Q)\le d(P,Z)+d(Z,Q)< r+r = d(P,Q),
\]
矛盾。因此 $U\cap V=\varnothing$。 
\end{dxtips}
\begin{homework}[7]
证明：$\mathbb{R}^2$ 是 $T_2$ 空间。
\end{homework}

\begin{homework}[5.3]
证明：离散空间 $X$ 是紧空间当且仅当 $X$ 是有限集。
\end{homework}

\section{十三周作业}
\begin{dxtips}
    紧空间加连续映射等于有极值，这种二级结论明天早上总结还有什么紧的闭子集是紧的
\end{dxtips}
\begin{homework}[5.5]
设 $f$ 是紧空间 $X$ 上的实值连续映射，且对任意 $x\in X$ 都有 $f(x)>0$。证明存在 $\varepsilon>0$，使对任意 $x\in X$ 都有 $f(x)>\varepsilon$。
\end{homework}

\begin{homework}[10]
$f:X\to Y$ 是连续满映射，若 $X$ 是紧空间且 $Y$ 是 $T_2$ 空间，证明 $f$ 是闭映射。
\end{homework}

\begin{homework}[11]
设 $\mathbb R$ 为实数直线，拓扑是通常拓扑，判断下列集合是否是紧集：
\begin{enumerate}
  \item $A=[1,2]$；
  \item $A=[1,2]\cup\{3\}$；
  \item $A=\mathbb{Q}$；
  \item $A=[3,4)$；
  \item $A=\mathbb{R}\setminus\{1/n:n\in\mathbb{N}\}$；
  \item $A=\mathbb{Z}$；
  \item $A=[1,3)\cup(2,4]$；
  \item $A=\{1/n:n\in\mathbb{N}\}$；
  \item $A=\{1/n:n\in\mathbb{N}\}\cup\{0\}$；
  \item $A=\{6-1/n:n\in\mathbb{N}\}\cup\{4\}$；
  \item $A=\{6-1/n:n\in\mathbb{N}\}\cup\{6\}$；
  \item $A=\{\pi,\pi+1\}$。
\end{enumerate}
\end{homework}
\begin{dxtips}
    和前面有到题很像，想证明A是闭集就证明A的补集是开集，A的补集中的x都要找到一个开邻域，先把包含x但是和A不相交的集族找出来，对应的是Y中不包含x的开集族，这个开集族覆盖了A找可数子覆盖，有限子覆盖的开集去包含x的找到对应的有限子覆盖，这个开集的交还是开集因为是有限交，然后这个交集就是x的开邻域因为和A没有交集，这个开邻域就在补集里面，知道有界闭集了，他肯定属于[-a，a]这个有界集这个条件是怎么都可以用的，然后因为是他的闭子空间所以他也紧
\end{dxtips}
\begin{homework}[4]
证明：集合 $A$ 是 $\mathbb R$ 中的紧集当且仅当 $A$ 是 $\mathbb R$ 中的有界闭集。
\end{homework}

\begin{homework}[12]
设 $\mathbb{R}$ 为实数直线，拓扑为通常拓扑，判断下列各集合是否为紧集：
\begin{enumerate}
  \item $A=(1,2)\times(3,4)$；
  \item $A=[1,2]\times\{3\}$；
  \item $A=\mathbb{Q}\times\mathbb{Q}$；
  \item $A=\{(x,y):x^{2}+y^{2}\le 4\}$；
  \item $A=\{(x,y):x+y=4\}$；
  \item $A=[1,2]\times[3,4]$；
  \item $A=\{(x,y):2\le x^{2}+y^{2}\le 4\}$。
\end{enumerate}
\end{homework}
\section{十四周作业}
\begin{homework}[13]
设
\[
X=\left\{\frac{1}{n}:n\in\mathbb{N}\right\}\cup I,
\]
其中 $I$ 为无理数集，$\mathbb{N}$ 为正整数集。

对任意 $n\in\mathbb{N}$，规定
\[
\mathcal{B}\!\left(\tfrac{1}{n}\right)=\bigl\{\{\tfrac{1}{n}\}\bigr\}
\]
为 $\tfrac{1}{n}$ 在 $X$ 中的开邻域基；

对任意 $x\in I$，规定
\[
\mathcal{B}(x)
=
\bigl\{\{\tfrac{1}{n}:n\ge m\}\cup\{x\}:m\in\mathbb{N}\bigr\}
\]
为 $x$ 在 $X$ 中的开邻域基。

证明下列结论：
\begin{enumerate}
\begin{dxtips}
    紧空间任意开覆盖都需要有限子覆盖，有限子覆盖只能覆盖有有限个1/n。第二题每一个x都需要有可数邻域基，题目中已经给了每一个x可数开邻域基了。
    第二可数要求有一个可数开邻域基，所有的开集都能用这个可数开邻域基中元素表述，如果有可数开邻域基，那么任意包含x的开集都有基中元素Bx属于Ux，可以等于就是基中总有比任何开集都小的集合因为基本来就是可以生成所有开集。
\end{dxtips}
  \item $X$ 不是紧空间；
  \item $X$ 是第一可数空间；
  \item $X$ 不是第二可数空间；
  \begin{dxtips}
      可数稠密子集就是B（1/n）因为任意开集只要有无理数就只能写成无理数并上1/n n大于m这个尾巴这个尾巴和D永远有交集就保证了稠密
  \end{dxtips}
  \item $X$ 是可分空间；
  这个要求只要x不等于y就有包含x和y的开集互不相交，但是因为包含x和y的开集都有尾巴也就是n大于m{1/n}所以总会有交集。
  \item $X$ 不是 $T_2$ 空间。
\end{enumerate}
\end{homework}
\begin{homework}[14]
若 $X$ 是紧的连通空间，$f:X\to\mathbb{R}$ 为连续映射
\end{homework}
\chapter{连续映射}
\begin{homework}[1]
证明：\(X\) 与 \(Y\) 是拓扑空间且 \(X\) 是离散拓扑空间。  
则任一映射 \(f:X\to Y\) 都是连续映射。
\end{homework}

\begin{homework}[2]
证明：\(X\) 与 \(Y\) 是拓扑空间且 \(Y\) 是平凡拓扑空间。  
则任一映射 \(f:X\to Y\) 都是连续映射。
\end{homework}

\begin{homework}[5]
设 \(\mathbb{R}\) 是实数直线，拓扑为通常拓扑。  
定义映射 \(f:\mathbb{R}\to\mathbb{R}\) 为
\[
f(x)=
\begin{cases}
2, & x\ge0,\\
5, & x<0.
\end{cases}
\]
证明：\(f\) 不是连续映射。
\end{homework}

\begin{homework}[6]
设 \(\mathbb{R}\) 是实数直线，拓扑为离散拓扑。  
定义映射 \(f:\mathbb{R}\to\mathbb{R}\) 为
\[
f(x)=
\begin{cases}
2, & x\ge0,\\
5, & x<0.
\end{cases}
\]
证明：\(f\) 是连续映射。
\end{homework}

\begin{homework}[7]
设 $X=\{a,b,c\}$，拓扑为 
\[
\mathcal{T}_X=\{X,\varnothing,\{a\},\{a,b\},\{a,c\}\},
\]
又设
\[
Y=\{a,b,c\},\quad 
\mathcal{T}_Y=\{Y,\varnothing,\{a\},\{a,b\},\{a,c\}\}.
\]
若映射 $f:X\to Y$ 满足 $f(x)=x$，问：$f$ 是否连续？
\end{homework}

\begin{homework}[8]
设 $X=\{a,b,c\}$，拓扑为 
\[
\mathcal{T}_X=\{X,\varnothing,\{a\},\{a,b\},\{a,c\}\},
\]
又设
\[
Y=\{a,b\},\quad 
\mathcal{T}_Y=\{Y,\varnothing,\{b\}\}.
\]
若映射 $f:X\to Y$ 满足 
\[
f(a)=b,\quad f(b)=b,\quad f(c)=a,
\]
问：$f$ 是否连续？
\end{homework}

\begin{homework}[9]
设 $X=\{a,b,c\}$，拓扑为 
\[
\mathcal{T}_X=\{X,\varnothing,\{a\},\{a,b\},\{a,c\}\},
\]
又设
\[
Y=\{a,b\},\quad 
\mathcal{T}_Y=\{Y,\varnothing,\{b\}\}.
\]
若映射 $f:X\to Y$ 满足 
\[
f(a)=a,\quad f(b)=b,\quad f(c)=a,
\]
问：$f$ 是否连续？
\end{homework}

\begin{homework}[12]
证明：映射 $f:X\to Y$ 连续，当且仅当 $f:X\to f(X)$ 连续。
\end{homework}

\begin{homework}[3]
证明：Lindelöf 拓扑空间的连续像是 Lindelöf 空间。
\end{homework}

\begin{homework}[4]
证明：可分拓扑空间的连续像是可分空间。
\end{homework}

\begin{homework}[14]
如果 $f:X\to Y$ 是连续映射，$A\subset X$，证明：$f|_A:A\to Y$ 也是连续映射。
\end{homework}

\begin{homework}[3.2]
若 $f_i:X\to\mathbb{R}$ 是连续映射，$i\in\{1,2\}$，证明：
$f_1+f_2$、$f_1-f_2$、$f_1f_2$ 及 $f_1/f_2$ 都是连续映射。
\end{homework}

\begin{homework}[3.6]
设 $f : X \to Y$ 是连续映射，证明：$f$ 是闭映射当且仅当对每个 $A \subset X$ 都有
\[
f(\overline{A}) = \overline{f(A)}.
\]
\end{homework}

\begin{homework}[3.3]
若 $f_n : X \to [a,b]$ 是连续映射，$n \in \mathbb{N}$，证明若
\[
f(x) = \sum_{n=1}^{\infty} \frac{f_n(x)}{2^n},
\]
则 $f : X \to [a,b]$ 是连续映射。
\end{homework}

\begin{homework}[2]
证明连通空间的连续像是连通空间。
\end{homework}

\begin{homework}[4.6]
设 $X$ 是连通空间，$f:X\to\mathbb R$ 连续。若 $x,y\in X$ 且 $f(x)f(y)<0$，证明存在 $z\in X$ 使得 $f(z)=0$。
\end{homework}

\begin{homework}[2]
证明紧子空间的像都是紧子空间。
\end{homework}

\begin{homework}[10]
$f:X\to Y$ 是连续满映射，若 $X$ 是紧空间且 $Y$ 是 $T_2$ 空间，证明 $f$ 是闭映射。
\end{homework}

\begin{homework}[5.5]
设 $f$ 是紧空间 $X$ 上的实值连续映射，且对任意 $x\in X$ 都有 $f(x)>0$。证明存在 $\varepsilon>0$，使对任意 $x\in X$ 都有 $f(x)>\varepsilon$。
\end{homework}

\begin{homework}[14]
若 $X$ 是紧的连通空间，$f:X\to\mathbb{R}$ 为连续映射
\end{homework}